\chapter{Análise Assintótica e Complexidade}
\label{cap:analise}

% TODO(P2): Apresentar análise matemática rigorosa de tempo e espaço com notações O, Omega e Theta:
%
% 1. Método Geral de Análise em Programação Dinâmica:
%    - Tempo Total = (Número total de subproblemas distintos) * (Custo de decisão/combinação por subproblema).
%    - Espaço Total = (Tamanho da tabela de armazenamento) + (Profundidade da pilha de recursão no Top-Down).
%
% 2. Prova Formal de Ineficiência do Corte de Hastes Ingênuo:
%    - Recorrência: T(0) = 1 e T(n) = 1 + \sum_{j=0}^{n-1} T(j).
%    - Demonstração por indução matemática de que T(n) = 2^n = \Theta(2^n).
%
% 3. Análise dos Algoritmos Otimizados:
%    - Fibonacci: Naive \Theta(\phi^n) vs. PD \Theta(n) tempo e \Theta(1) espaço.
%    - Corte de Hastes: Bottom-Up \Theta(n^2) tempo e \Theta(n) espaço.
%    - Cadeia de Matrizes: \Theta(n^3) tempo e \Theta(n^2) espaço (laço triplo sobre subcadeias e pontos de partição).
%    - LCS: \Theta(m \cdot n) tempo e \Theta(m \cdot n) espaço para tabela completa, reduzido a \Theta(\min(m, n)) espaço com duas linhas.
%
% 4. Tabela-Resumo Oficial de Complexidades:
%    - Construir tabela com colunas: Problema | Abordagem | Tempo (Pior/Médio) | Espaço | Observações.
%    - Esta tabela será espelhada diretamente no Cheat Sheet oficial.

% [Texto e deduções matemáticas a cargo de P2 - ver docs/tarefas/P2_relatorio_secoes_3_4.md]
